\documentclass[10pt,a4paper]{amsart}
\usepackage[margin=19mm]{geometry}
\usepackage{amsmath,amssymb,mathtools}
\usepackage[scr=rsfs,cal=euler]{mathalfa}
\usepackage{xfrac}
\usepackage[unicode,hidelinks,bookmarks=false]{hyperref}
\newcommand{\ie}{i.e.\ }
\newcommand{\cf}{cf.\ }
\newcommand{\resp}{respectively\ }
\newcommand{\Subsection}{\subsection}
\providecommand{\nobreakdash}{\nobreak\hbox{-}}
\setlength{\emergencystretch}{2em}
\multlinegap0pt
\usepackage{amssymb}
\usepackage{mathtools}
\usepackage{float}
\usepackage{tikz}
\usepackage[shortlabels]{enumitem}
\newtheorem{theorem}{Theorem}
\title{Oscillatory motion to collision and infinity in the Earth-Moon restricted three body problem}
\author{Aleksander Pasiut}
\date{Source: arXiv:2608.05400v3 (CC BY 4.0)}
\begin{document}
\maketitle

\noindent\emph{Source and licence.} Oscillatory motion to collision and infinity in the Earth-Moon restricted three body problem. Version arXiv:2608.05400v3 (CC BY 4.0), distributed by arXiv under the Creative Commons Attribution 4.0 International licence, http://creativecommons.org/licenses/by/4.0/, as recorded in the arXiv licence metadata for this exact version and shown on the abstract page https://arxiv.org/abs/2608.05400v3. Full licence notice: \texttt{licenses/arxiv-2608.05400-CC-BY-4.0.txt}.\par\smallskip

\noindent\textit{Editorial scope.} Counted unit: the article's theorem thm:exp\_trick (source0.tex lines 1281--1308). The article's complete original proof of it is delivered with it (source0.tex lines 1309--1328, one \texttt{proof} environment of 20 lines). No mathematics is written, repaired or re-derived: the statement and the proof are the article's own lines, in the article's own notation, and no retained line is substituted or re-flowed.  No further source-local context is needed: every cross-reference of every retained line resolves inside the two delivered spans.  Nothing was omitted from any retained span, so the package carries no omission record.  No retained line contains a citation, so the wrapper carries no bibliography and no bibliography entry.  No retained line contains a figure environment, an image-inclusion command or drawing code, so no image is delivered and the package declares no asset dependency.  Editorial formatting only, in addition: an AMS \texttt{amsart} preamble that restates the article's own declarations and macros used by the retained lines, namely the environments \texttt{theorem} on separate counters, together with the standard packages those lines need.  The retained environments print in this document as Theorem 1 in order of appearance, whereas the article prints the same items with the numbering of its own full document.  The complete native source of the article as distributed in the arXiv e-print for this exact version is bundled unchanged as \texttt{sources/206922/source0.tex}.
\begin{theorem}
    \label{thm:exp_trick}
    Let $n$ and $m$ be the positive integer values. We consider an ordinary differential equation:
    \begin{equation}
        \label{eq:ode_initial}
        \begin{cases}
            \dot{X} &= \mathcal{B}(X) \mathcal{A}(X,Y) X, \\
            \dot{Y} &= \mathcal{F}(X,Y), \\
        \end{cases}
    \end{equation}
    with the initial condition $X(0) = X_0 \in \mathbb{R}^n$, $Y(0) = Y_0 \in \mathbb{R}^m$, where:
    \[
        \mathcal{B} : \mathbb{R}^n \to \mathbb{R}, \quad
        \mathcal{A} : \mathbb{R}^{n+m} \to \mathbb{R}^{n \times n}, \quad
        \mathcal{F} : \mathbb{R}^{n+m} \to \mathbb{R}^m.
    \]
    The solution $X(t)$ of the equation (\ref{eq:ode_initial}) satisfies the equation:
    \begin{equation}
        \label{eq:solution_x}
        X(t) - X(0) = \mathcal{B}(X_0) U(t),
    \end{equation}
    where the functions $C(t)$ and $U(t)$ are the solution of the system of the ordinary differential equations:
    \begin{align}
        \dot{C} &= \big( D \mathcal{B}(X) \mathcal{A}(X,Y) X \big) C, \label{eq:ode_c} \\
        \dot{U} &= C \cdot \mathcal{A}(X,Y) \big( X_0 + \mathcal{B}(X_0) U \big), \label{eq:ode_gamma}
    \end{align}
    with the initial condition $C(0) = 1$ and $U(0) = 0 \in \mathbb{R}^n$.
\end{theorem}

\begin{proof}
    We notice that the function $t \mapsto \mathcal{B}(X(t))$ is the solution of the differential equation:
    \begin{equation}
        \label{eq:ode_b}
        \frac{d}{dt} \mathcal{B}(X(t)) = \big( D \mathcal{B}( X ) \mathcal{A}(X,Y) X \big) \mathcal{B}(X(t)),
    \end{equation}
    which means that $\mathcal{B}(X(t))$ can be expressed as
    \begin{equation}
        \label{eq:b_solution}
        \mathcal{B}(X(t)) = \mathcal{B}(X_0) C.
    \end{equation}

    We differentiate the equation (\ref{eq:solution_x}) with respect to $t$ and we substitute (\ref{eq:ode_gamma}) into the result of the differentiation, which leads us to
    \begin{equation}
        \label{eq:ode_x_2}
        \dot{X}(t) = \mathcal{B}(X_0) C \cdot \mathcal{A}(X,Y) \big( X_0 + \mathcal{B}(X_0) U(t) \big).
    \end{equation}

    We notice that the equation above is equivalent to the original one (due to the relationships (\ref{eq:b_solution}) and (\ref{eq:solution_x})), which ends our proof.
\end{proof}

\end{document}
